% !TEX program = xelatex
\documentclass{英语试卷}
\begin{document}
\examtitle{绝密★启用前 鄂南高中 黄冈中学 黄石二中 荆州中学 龙泉中学}{英语}
\par 绝密★启用前 鄂南高中 黄冈中学 黄石二中 荆州中学 龙泉中学\par
\par 武汉二中 孝感高中 襄阳四中 襄阳五中 宜昌一中 夷陵中学\par
\par 英语试题\par
\par 命题学校：夷陵中学 命题人：黄家清 向晓晓 罗芬 审题学校：龙泉中学\par
\parttitle{第一部分 听力（共两节，满分 30 分）}
\par 做题时，先将答案标在试卷上。录音内容结束后，你将有两分钟的时间将试卷上的答案转涂到答题卡上。\par
\sectitle{第一节（共5小题；每小题1.5分，满分7.5分）}
\par 听下面 5 段对话。每段对话后有一个小题，从题中所给的 A、B、C 三个选项中选出最佳选项，并标在试卷的相应位置。听完每段对话后，你都有 10 秒钟的时间来回答有关小题和阅读下一小题。每段对话仅读一遍。\par
\question{1. What do we know about Tom?}
\begin{choices}(2)
  \item He's just finished a piano lesson. 3. He loves playing the guitar.
  \item He plays drums in a band.
\end{choices}
\question{2. Where are the speakers probably?}
\begin{choices}(3)
  \item At a café.
  \item At a university.
  \item At an animal shelter.
\end{choices}
\question{3. How can the man's recipe be described?}
\begin{choices}(3)
  \item It's a reward.
  \item It's a success.
  \item It's a problem.
\end{choices}
\question{4. What do the two speakers have in common?}
\begin{choices}(3)
  \item They enjoy reading.
  \item They are in the library.
  \item They have just finished a book.
\end{choices}
\question{5. What is the woman's duty in the event?}
\begin{choices}(3)
  \item To make food.
  \item To play in the band.
  \item To give directions.
\end{choices}
\sectitle{第二节（共15小题；每小题1.5分，满分22.5分）}
\par 听下面 5 段对话或独白。每段对话或独白后有几个小题，从题中所给的 A、B、C 三个选项中选出最佳选项，并标在试卷的相应位置。听每段对话或独白前，你将有时间阅读各个小题，每小题 5 秒钟；听完后，各小题将给出 5 秒钟的作答时间。每段对话或独白读两遍。\par
\question{6. Why was it difficult to rescue the cat?}
\begin{choices}(3)
  \item It was too scared to move.
  \item It was hurt by the branches.
  \item It was in the top of a tree.
\end{choices}
\question{7. What did the man do to rescue the cat?}
\begin{choices}(3)
  \item He got a neighbor to help him.
  \item He borrowed a ladder.
  \item He climbed the tree.
\end{choices}
\question{8. What is the woman doing?}
\begin{choices}(3)
  \item Conducting an interview.
  \item Reporting some news.
  \item Organizing an event.
\end{choices}
\question{9. What is Charles Martin's family mainly worried about?}
\begin{choices}(3)
  \item His health to become worse.
  \item His efforts not to be recognized.
  \item His neglect of his studies or family.
\end{choices}
\question{10. How does Charles Martin plan to remove his family's concerns?}
\begin{choices}(3)
  \item By talking to his teacher.
  \item By involving his family.
  \item By quitting the project.
\end{choices}
\par 听第 8 段材料，回答第 11 至 13 题。\par
\question{11. What are the speakers?}
\begin{choices}(3)
  \item A player and a coach.
  \item Football fans.
  \item TV sports reporters.
\end{choices}
\question{12. What causes the current situation of the game according to the man?}
\begin{choices}(3)
  \item Bad defense.
  \item Poor passing.
  \item Certain players.
\end{choices}
\question{13. What does the woman think will happen at the end of the game?}
\begin{choices}(3)
  \item The home team will win.
  \item The visiting team will win
  \item The game will end in a draw.
\end{choices}
\par 听第 9 段材料，回答第 14 至 17 题，\par
\question{14. What is the main topic of the conversation?}
\begin{choices}(3)
  \item The working life.
  \item College experience.
  \item Hobbies.
\end{choices}
\question{15. How long has the man been at university?}
\begin{choices}(3)
  \item About four months.
  \item About six months.
  \item About nine months.
\end{choices}
\question{16. What is the biggest issue for the man?}
\begin{choices}(3)
  \item The stress of exams.
  \item The shortage of money.
  \item The loads of assignments.
\end{choices}
\question{17. When does the man find time to relax?}
\begin{choices}(3)
  \item On weekends.
  \item Right after exams.
  \item After part-time work.
\end{choices}
\par 听第 10 段材料，回答第 18 至 20 题。\par
\question{18. How did the speaker become a photographer in demand?}
\begin{choices}(3)
  \item By teaching in workshops.
  \item By starting a business.
  \item By taking odd jobs.
\end{choices}
\question{19. What challenge did the speaker face?}
\begin{choices}(3)
  \item The stress of taking perfect pictures.
  \item Decreased interest in photography.
  \item Constantly changing confidence.
\end{choices}
\question{20. What is included in the speaker's future plans for his career?}
\begin{choices}(3)
  \item Taking photography trips abroad.
  \item Expanding his photography business.
  \item Learning more forms of photography.
\end{choices}
\parttitle{第二部分 阅读理解（共两节，满分 50 分）}
\sectitle{第一节（共15小题；每小题2.5分，满分37.5分）}
\par 阅读下列短文，从每题所给的 A、B、C 和 D 四个选项中，选出最佳答案，\par
\passagelabel{A}
\par With the growing popularity of the TV series Blossoms Shanghai, key scenes presenting city areas in the show have become the hottest tourist attractions in Shanghai of late.\par
\par The narrator in the TV series Blossoms Shanghai points out: "In 1993, the 755-meter-longHuanghe Road had over 100 restaurants, attracting countless wealthy diners each night." Starting from Huanghe Road and crossing Nanjing Road to the Bund, you'll pass by many scenic spots often filmed in the series:\par
\par · Wing On Department Store\par
\par The prototype(原型) of Hu Lian Department Store in the TV series was Hua Lian Departent Store between 1988 and 2006. It later changed its name to Wing On Department Store and was one of the four major original department stores in Shanghai.\par
\par Address: No 635 East Nanjing Road, Huangpu district.\par
\par · Fairmont Peace Hotel: Peace Hotel on East Nanjing Road.\par
\par It's a major scene in the TV series where the hero A Bao, played by actor Hu Ge, resides. The hotel first opened its doors on Aug 1, 1929, and was originally named Cathay Hotel by its owner, Sir Victor Sassoon, a businessman and hotelier.\par
\par Address: No 20 East Nanjing Road, Huangpu district.\par
\par · No 27 The Bund\par
\par As one of the important buildings in the Bund architectural complex and a key cultural relic protected unit in China, the building used to be owned by Jardine Matheson and later served as an office building for the Shanghai Foreign Trade Bureau for many years. In the TV series, the location is called the Foreign Trade Building.\par
\par Address: No 27 East Zhongshan Road, Huangpu district.\par
\question{21. According to the narrator in the Blossoms Shanghai, what can we know about the Huanghe Road?}
\begin{choices}(2)
  \item It is where the hero A Bao resides.
  \item It has the most restaurants in Shanghai.
  \item It appealed to many rich eaters every night.
  \item It is the site of the Wing On Department Store.
\end{choices}
\question{22. Which of the following buildings was initially named Cathay Hotel?}
\begin{choices}(2)
  \item No 27 The Bund.
  \item Fairmont Peace Hotel.
  \item Hu Lian Department Store.
  \item Wing On Department Store.
\end{choices}
\question{23. If you are interested in architecture, which place would you like to visit?}
\begin{choices}(2)
  \item No 755 Huanghe Road, Huangpu district.
  \item No 20 East Nanjing Road, Huangpu district.
  \item No 635 East Nanjing Road, Huangpu district.
  \item No 27 East Zhongshan Road, Huangpu district.
\end{choices}
\passagelabel{B}
\par Late one night in July 2020, Halli was walking around Reykjavik's main shopping street with his family. They had just moved back to Iceland from the United States and were struggling with jet lag(时差). When his three-year-old son wanted a drink from a corner store, Halli faced a problem: a20-centimetre step blocked his wheelchair from entering the store.\par
\par This was not new to Halli, who has used a wheelchair since age 25 due to a muscle disease. As he waited outside, he thought about how often he had been excluded from places like restaurants and stores because they were not accessible.\par
\par Halli, a successful digital designer, had lived in many cities worldwide and noticed how differently they approached accessibility. After selling his company, Ueno, to Twitter(now X), he decided to use his resources to make Iceland more wheelchair-friendly.\par
\par In 2021, he launched Ramp (斜坡) Up Reykjavik, a non-profit project aiming to build 100 permanent ramps in the city's downtown within a year. Unlike temporary ramps in other cities, these were designed to fit perfectly with the buildings, making them look like they had always been there.\par
\par With government funding and sponsors, the project went beyond its goal and expanded nationwide. By 2026, the team plans to build 1,500 ramps across Iceland, costing about 400 million Icelandic krona.\par
\par Iceland's president praised the project, saying, "We want our society to be inclusive. Halli has shown us how easy it is to make a difference when we work together."\par
\par In just three years, Halli has become a local hero, known as a "kind tech giant". He hopes his work inspires others. Recently, a woman built a ramp at her home so her wheelchair-using friend could visit easily. "Equal access is not a reality yet," he says, "but change starts with one person."\par
\question{24. What inspired Halli to start the 'Ramp Up Reykjavik' project?}
\begin{choices}(3)
  \item The lack of wheelchair access. 3. His son's difficulty in finding a drink.
  \item A suggestion from Iceland's president.
  \item His desire to become famous in Iceland.
\end{choices}
\question{25. How are the ramps built by Ramp Up Reykjavik' different from those in other cities?}
\begin{choices}(2)
  \item They are cheap to construct.
  \item They are only built in downtown.
  \item They are temporary and movable.
  \item They look fitted with the buildings forever.
\end{choices}
\question{26. Which of the following best describes Halli?}
\begin{choices}(2)
  \item Generous and patient.
  \item Positive and energetic.
  \item Smart and hardworking.
  \item Inspiring and pioneering.
\end{choices}
\question{27. What dose Halli expect people to do about equal access?}
\begin{choices}(2)
  \item Go ahead to start from small.
  \item Build ramps in every city
  \item Appeal to government for funds.
  \item Raise awareness through media.
\end{choices}
\passagelabel{C}
\par When someone proposes a false claim, what's the best way to change their mind? A recent paper suggests that immediately negating the claim with evidence isn't especially effective. Instead, "bypassing" the false claim with positive counter(反)-claims about the topic might be a better strategy.\par
\par The classic response to someone presenting misinformation is to present counter-evidence. As the researchers at the Annenberg Public Policy Center (APPC) put it, "The gold standard for tackling misinformation is usually a correction that actually contradicts the misinformation." Suppose, for example, that your dad passes back the roast potatoes over a family dinner, saying, "Are these GMOs(转基因产品)? Nah, don’t want any more physical reactions." You might then say, "Actually, the health bodies of all the G7 agree that GMOs do not cause reactions." But that likely won't work. The APPC team conducted six experiments where they compared actually correcting someone with "bypassing" and found that using evidence and facts to disprove someone tends to be far less successful at changing attitudes.\par
\par So, what is bypassing and why is it seemingly more effective? Bypassing is where, instead of providing some negative disapproval of a claim, you instead provide positive counter-claims about the topic. For example, if someone tells you that 5G kills birds, tell them that 5G adds billions to the world economy. Bypassing means not meeting your misinformed opponent head-on. This is not a battle of arguments or the studies. You are pushing an argument into a different position. It's using someone's movements and logic against them.\par
\par There are limitations to what the team at APPC has concluded. Their study into bypassing was largely about attitudes to policies, not about belief change or belief correction. Someone might still believe that 5G kills birds or that GMOs cause reactions, but admit that those are necessary bad things, compared to the positive outcomes you present them with.\par
\par So, bypassing isn't the only tool available, and it might not even be the best if your concern is belief change. Changing someone's belief is a complicated psychological topic, and there is definitely no one-size-fits-all strategy. When it comes to policy decisions and willingness to accept a certain technology or medicine, bypassing seems like a valuable tool to add to your kit.\par
\question{28. According to the APPC research, how does "bypassing" differ from traditional corrections?}
\begin{choices}(2)
  \item It uses stronger evidence.
  \item It adopts emotional appeals.
  \item It offers positive arguments.
  \item It focuses on unrelated topics.
\end{choices}
\question{29. Which can be used as an example of "bypassing"?}
\begin{choices}(2)
  \item GMOs are safe to eat because scientists say so
  \item GMOs can produce higher output with less cost.
  \item 5G kill birds, but their benefits outweigh the risks.
  \item 5G doesn't kill birds, for they show no connection.
\end{choices}
\question{30. What do we know about the "bypassing" tool according to the text?}
\begin{choices}(2)
  \item It is the most reliable way.
  \item It may help accept policies.
  \item It appears frequently in debates.
  \item It is effective in belief correction.
\end{choices}
\question{31. Which of the following is the best title for the text?}
\begin{choices}(2)
  \item Bypassing: A Better Way to Shift Attitudes?
  \item Changing one's Mind: Bypass Indirectly!
  \item Challenging Misinformation: Bypass it!
  \item Negating the claim: does it work?
\end{choices}
\passagelabel{D}
\par People frequently employ AI to craft emails, plan travel, get financial advice, summarize texts, and prepare for job interviews. According to a 2024 Pew Research survey, roughly half of Americans said they use AI at least several times a week, while a recent survey found that virtually all Americans use products that contain AI (even though two-thirds don't realize it).\par
\par Companies offering AI assistants sell them as productivity boosters. They argue that assigning them routine mental tasks frees up users' time and cognitive (认知的) resources, which they can then spend on other, more creative, and innovative pursuits. The idea makes natural sense and has backing in the scientific literature. According to cognitive load theory, the human cognitive system has limited capacity. So reducing cognitive load can enhance learning and performance.\par
\par Habitually offloading cognitive tasks to AI could backfire, however. As AI has grown more commonplace in everyday life, psychologists theorize that it reduces users' engagement in deep, reflective thinking, causing their critical thinking skills to decline over time.\par
\par Professor Dr. Michael Gerlich, Head of the Center for Strategic Corporate Foresight and Sustainability at the Swiss Business School, is one of the researchers studying this risk. "If individuals use the cognitive resources freed up by AI for innovative tasks, the promise holds." he told Big Think.\par
\par "However, my research and related studies suggest that many users channel these resources into pass: consumption, driven by AI-enhanced content selection. This trend matches findings on digital dependence, where the convenience of AI encourages a feedback cycle that prioritizes entertainment over critical engagement."\par
\par In other words, when AI frees up users' cognitive resources, they typically don't use their extra time and brain power to problem-solve or create. Rather, they tune out by watching Netflix or browsing social media — content served up by AI algorithms (算法). Gerlich found a very strong negative correlation between subjects' use of AI tools and their critical thinking skills. The higher their usage, the lower their skills. Younger participants tended to be more dependent on AI tools compared to older participants. Education was associated with greater critical thinking skills and reduced AI's negative effect.\par
\question{32. What is the main argument of companies offering AI assistants?}
\begin{choices}(2)
  \item They improve efficiency and output.
  \item They help users reduce screen time.
  \item They expand users' cognitive abilities.
  \item They enhance creativity and innovation.
\end{choices}
\question{33. Which can replace the underlined word "backfire" in paragraph 3?}
\begin{choices}(2)
  \item Result in a positive effect.
  \item Cause negative consequences.
  \item Achieve an intended outcome.
  \item Prove the existing theory wrong.
\end{choices}
\question{34. What does Dr. Gerlich's research suggest about the use of AI?}
\begin{choices}(2)
  \item Freed-up resources fue innovation.
  \item AI is mainly for educational purposes
  \item Older users are more likely to rely on AI.
  \item Freed-up resources go to mindless activities.
\end{choices}
\question{35. Which can describe the relationship between AI usage and critical thinking?}
\begin{choices}(2)
  \item Irrelevant.
  \item Complicated.
  \item Disadvantageous.
  \item Positive.
\end{choices}
\sectitle{第二节（共5小题；每小题2.5分，满分12.5分）}
\par 根据短文内容，从短文后的选项中选出能填入空白处的最佳选项。选项中有两项为多余选项，\par
\par We've all experienced the frustration of forgetting what we studied yesterday. But memory isn't just about natural talent. It's a skill we can improve with smart strategies.\par
\par Our brains naturally screen information. We forget 70\% of new knowledge within 24 hours if we don't review it. Your brain serves as a library. Whenever you revisit information, you're placing it on a more accessible shelf. A student who reviews math formulas(公式) right after class, then before bed, and once more after three days will find them easier to recall during exams. \blanknum{36}\par
\par \blanknum{37} When reciting words like "photosynthesis", simply repeating the spelling often fails quickly. Instead, picture sunlight hitting green leaves while reading a biology textbook. This combines text, imagery, and real-world meaning, activating multiple brain areas at once. Similarly, linking "pineapple" to "pine tree with apples" creates a funny mental image hard to forget for long.\par
\par Sleep and calm minds are secret weapons for memory. During deep sleep, the brain reorganizes daytime learning like a librarian sorting books. Research shows students with 7-8 hours of sleep remember \blanknum{40}\% more history facts than those who stay up late. \blanknum{38} Taking slow breaths can protect your mental "library".\par
\par Active practice beats passive reading. Instead of mechanically rereading notes, transform headings into self-test questions: "\blanknum{39}" This forces your brain to search information, strengthening memory pathways. Explaining concepts to friends, like teaching a classmate how to solve equations, uses speech and logic to deepen understanding.\par
\par \blanknum{40} By spacing out reviews, creating vivid mental connections, and caring for your brain's needs, you can turn temporary knowledge into lifelong wisdom — one small, steady step at a time.\par
\begin{choices}(1)
  \item How do I use this grammar rule in writing?
  \item This "spaced repetition" makes memories stick.
  \item Forgetting is a common human condition without exception.
  \item Context turns weak and uncertain memories into permanent ones
  \item Memory mastery is not about studying longer but studying smarter.
  \item Improving memory is more focused on holding on to what really counts.
  \item Feeling stressed, however, floods the brain with chemicals that damage memory.
\end{choices}
\parttitle{第三部分 语言知识运用（共两节，满分 30 分）}
\sectitle{第一节（共15小题；每小题1分，满分15分）}
\par 阅读下面短文，从短文后各题所给的 A、B、C 和 D 四个选项中，选出可以填入空白处的最佳选项。\par
\par Bill is a Navy sailor and retired oilfield engineer. His hands are \blanknum{41} and thick. Pulling on ropes as a sail or off the coast of Maine gnarled(扭曲) a few of his fingers. The others were \blanknum{42} playing catcher(接球手) in South Africa. He has a story for every bone break, and for why he \blanknum{43} a bit to one side.\par
\par After returning to America, Bill found himself with \blanknum{44} time on hand. He started watching the local high school baseball team play. Baseball is a game of skill, focus and patience. So it is the same with \blanknum{45} the young men. Bill simply sat in the \blanknum{46} and carefully watched them. Due to his former experience, it \blanknum{47} Bill to watch the boys develop their skills.\par
\par Soon, the boys noticed Bill. He \blanknum{48} his elbows (胳膊肘) on his knees, almost bending in his \blanknum{49} catcher's squat(蹲). After a good play, Bill would flash a thumbs-up. An error would get a smile, meaning, "You'll \blanknum{50} another chance."\par
\par It's rare to see a bond that \blanknum{51} generation divides — Bill in his late 70s and the boys just teenagers. But soon they were chatting after the games, Bill mostly talking about dusty adventures in \blanknum{52} of new oil fields and the boys asking about playing catcher. Sometimes, they would have sodas together, sharing laughs and \blanknum{53} conversations.\par
\par Bill has learned one important \blanknum{54} from this: just being there for others is often enough. His \blanknum{55} encouraged the boys to push themselves harder.\par
\cz{41}{soft}{smooth}{delicate}{rough}
\cz{42}{lost}{hurt}{ruined}{cut}
\cz{43}{goes}{runs}{tips}{falls}
\cz{44}{busy}{limited}{free}{valuable}
\cz{45}{blaming}{guiding}{accepting}{punishing}
\cz{46}{stands}{shelters}{woods}{stores}
\cz{47}{moved}{upset}{amazed}{suited}
\cz{48}{hidden}{rested}{raised}{stuck}
\cz{49}{brilliant}{strange}{familiar}{elegant}
\cz{50}{get}{lose}{miss}{take}
\cz{51}{connects}{links}{bridges}{creates}
\cz{52}{view}{search}{favour}{pursuit}
\cz{53}{dull}{mild}{private}{light}
\cz{54}{lesson}{skill}{principle}{fact}
\cz{55}{preference}{existence}{confidence}{presence}
\sectitle{第二节（共10小题；每小题1.5分，满分15分）}
\par 阅读下面短文，在空白处填入 1 个适当的单词或括号内单词的正确形式。\par
\par With just over a week to go until the opening of the \blanknum{56} (nine) Asian Winter Games, the streets of Harbin are alive with both the festive spirit of the Spring Festival in the Year of the Snake and growing excitement for \blanknum{57} upcoming sporting event.\par
\par At the Harbin Ice-Snow World, a theme park dedicated to the Asian Winter Games, visitors \blanknum{58} (treat) with a splendid drone(无人机) show celebrating the Chinese New Year. The park features a variety of ice sculptures with festive themes, including the Chinese character Fu \blanknum{59} (carve) in ice.\par
\par The world's \blanknum{60} (large) ice park also showcases ice sculptures of landmarks from various member countries and regions of the Olympic Council of Asia and snow sculptures \blanknum{61} (picture) different winter sports.\par
\par Harbin has seen boosting \blanknum{62} (enthusiastic) for ice and snow activities ahead of the Games. The entrances of shops are decorated with red couplets, creating the festive atmosphere. Preparations for the Asian Winter Games continued \blanknum{63} the Chinese New Year's Eve, when the Games' media village officially opened.\par
\par The New Paris Hotel, located on Zhongshan Road in Harbin, is one of the designated media reception hotels. \blanknum{64} (accommodate) international journalists, the hotel has arranged for volunteers expert in multiple languages. A media service desk has been set up in the hotel hall, \blanknum{65} staff are to make final arrangements for check-ins. Final preparations for the Games are in full swing during the Spring Festival.\par
\par ```txt Dear Ryan, I'm writing to share with you a popular Chinese movie I saw during this Spring Festival. \blanknum{} Yours, Li Hua ```\par
\parttitle{第四部分 写作（共两节，满分 40 分）}
\sectitle{第一节（满分15分）}
\par 假如你是李华，春节期间你看了一部好评如潮的国产电影。请你给英国朋友 David 写一封邮件分享这次经历，内容包括：(1) 你观看的电影；(2)你的感想， 注意：\par
\par (1) 写作词数应为 80 左右;\par
\par (2) 可适当增加细节，以使行文连贯，\par
\sectitle{第二节（满分25分）}
\par 阅读下面材料，根据其内容和所给段落开头语续写两段，使之构成一篇完整的短文。\par
\par The iron gates of Sunshine Nursing Home got opened as 16-year-old Tim stepped inside, his sneakers stepping on the autumn leaves on the ground. The hall was filled with the smell of air freshener. Elderly residents sat nodding in wheelchairs or played chess with trembling hands, lost in their own worlds. His volunteering here was a part of his school requirements.\par
\par Every Saturday, Tim found himself unwillingly volunteering at Sunshine Nursing Home. His tasks included assisting Mr. Johnson in the woodworking room and reading newspapers to Mrs. Clark in her room. Although he was initially forced, he began to find a sense of purpose in his work.\par
\par As he wandered through the nursing home, Tim found Mr. Johnson's woodworking room lay away behind the laundry area. The walls were decorated with half-carved birds and a sign that read, "Measure twice, cut once." The 78-year-old craftsman, with tough hands, guided Tim's hands as they sanded(打磨) wood boards together. When Tim accidentally hurt a table leg with his iron tool, Mr. Johnson said, "Wood remembers every mistake. Wood never lies. Let your hands feel its truth."\par
\par On the other side of the building, Mrs. Clark's sunlit room was a clear contrast. Overflowing with many books, it was a shelter for the former teacher, who was legally blind. Whenever Tim read newspaper articles to her, she would tap her pearl-ringed finger gently. "Young man," she encouraged him, putting her head to one side so her cat-eye glasses caught the light, "Words dance when you give them breath, dear. Even stock market reports deserve passion!" Her laughter, like wind chimes(风铃), covered up the loneliness Tim sometimes saw when her fingers felt over her family photos on the piano.\par
\par As the weeks passed, Tim learned to carve maple leaves under Mr. Johnson's guidance and practiced reading with more expression to Mrs. Clark. One October afternoon, both the elders fell silent when they heard about the nursing home's upcoming talent show. Tim noticed Mrs. Clark's trembling hands catching an old violin case under her armchair, meaning something she intended to do.\par
\par 注意：(1) 续写词数应为 150 左右；(2) 请按如下格式在答题卡的相应位置作答。\par
\writingpara{Paragraph 2: Tim's initial stutter(口吃) turned into a smooth story, his voice accompanied by Mrs. Clark's music.}
\writinglines{8}
\end{document}
